\begin{tikzpicture}[x=1cm,y=1cm,font=\small]
\draw[gray!45,rounded corners=3pt,fill=white] (0,0) rectangle (13.8,8.2);
\node[anchor=west,font=\small\bfseries] at (.4,7.7) {斜率相反，还要检查射线朝向，才能判断内角平分线};
\foreach \xx/\kind in {3.5/1,10.3/2}{
\begin{scope}[shift={(\xx,4.5)}]
\draw[cnmirror] (0,-2)--(0,2.1);
\ifnum\kind=1
 \fill[green!45!black,opacity=.08] (0,0)--(-1,-2)--(1,-2)--cycle;
 \draw[cnsec,-Stealth] (0,0)--(-1,-2);
 \draw[cnsec,-Stealth] (0,0)--(1,-2);
 \node[left] at (-1,-2) {$PM$}; \node[right] at (1,-2) {$PN$};
 \node[align=center] at (0,-3) {$m=2$，射线分居轴两侧\\竖直轴平分内角};
\else
 \fill[orange!85!black,opacity=.08] (0,0)--(-2,-1)--(-2,1)--cycle;
 \draw[cnsec,-Stealth] (0,0)--(-2,-1);
 \draw[cnsec,-Stealth] (0,0)--(-2,1);
 \node[left] at (-2,-1) {$PM$}; \node[left] at (-2,1) {$PN$};
 \draw[gray] (0,0)--(-2.3,0);
 \node[align=center] at (0,-3) {$m=1/2$，射线位于轴同侧\\竖直轴平分外角};
\fi
\fill (0,0) circle(1.8pt); \node at (.5,.35) {$P$};
\end{scope}}
\node at (6.9,.5) {仅画例中两条边射线的方向，不表示交点到 $P$ 的距离。};
\end{tikzpicture}
